% =====================================================================
\chapter{Measurement Framework}\label{ch:framework}
% =====================================================================

\begin{plan}[Chapter plan --- Measurement framework]
\pages{9 pages.} \textbf{Job:} define every quantity once, precisely, with the
assumptions that make it valid. This is the mathematical core of the thesis.
The definitions, equations and propositions below are \emph{planned final
content}; write the connecting prose around them. Keep the numbering of the
proposal's Eqs.~3.1--3.6 recognisable: the proposal's $P(\text{yes})$ is
\cref{eq:binary}, $b_i$ is \cref{eq:bias}, $\ASIint$ is \cref{eq:asi},
$\dm$ is \cref{eq:deltam}, $\RR$ is \cref{eq:rr}.
\citeplan{holtzman2021surface,zhao2021calibrate,baumgartner2001response,weijters2013reversed,ye2025justice,ernst2004permutation,efron1993bootstrap,field2007bootstrapping,zhang2019rmsnorm}
\end{plan}

% ---------------------------------------------------------------------
\section{Setting and notation}\label{sec:framework-setting}
\begin{plan}
Half a page. A measurement is one forward pass of model $m$ on prompt $x$ and
image $v$; generation is never run. Everything is a function of the
final-position distribution $p(\cdot\mid x,v)$ in \cref{eq:nexttoken}. Point
to the notation table (\cref{tab:notation}).
\end{plan}

% ---------------------------------------------------------------------
\section{The pooled binary probe and captured mass}\label{sec:framework-probe}
\begin{definition}[Pooled masses, captured mass and $\pyes$]\label{def:pooled}
With $\Yset=\{\text{yes},\text{Yes},\text{YES}\}$ and
$\Nset=\{\text{no},\text{No},\text{NO}\}$,
\begin{align}
  \pi_{\Yset} &= \sum_{y\in\Yset} p(y\mid x,v), \qquad
  \pi_{\Nset} = \sum_{n\in\Nset} p(n\mid x,v), \label{eq:pooled}\\
  \capm &= \pi_{\Yset} + \pi_{\Nset}, \label{eq:capm}\\
  \pyes &= \frac{\pi_{\Yset}}{\capm}, \qquad
  \logit(\pyes) = \operatorname*{LSE}_{y\in\Yset} z_y - \operatorname*{LSE}_{n\in\Nset} z_n .
  \label{eq:pyes}
\end{align}
\end{definition}

\begin{remark}[The proposal's two-token softmax is a special case]
With $\Yset=\{\text{yes}\}$ and $\Nset=\{\text{no}\}$, \cref{eq:pyes} reduces to
\begin{equation}
  \pyes = \frac{\exp(z_{\text{yes}})}{\exp(z_{\text{yes}})+\exp(z_{\text{no}})},
  \label{eq:binary}
\end{equation}
the proposal's Eq.~3.1. The normalising constant of \cref{eq:nexttoken} cancels,
so \cref{eq:binary} is the same whether the two tokens hold 99\% or 0.09\% of
the distribution. That is why $\capm$ must be reported alongside every
$\pyes$: it is the only quantity that says how much of the model's behaviour
$\pyes$ describes.
\end{remark}

\begin{plan}
\begin{itemize}
  \item Justify pooling with surface-form competition \cite{holtzman2021surface}
        and give the measured case: lowercase-only gave $\capm=0.09\%$ and
        overstated $\pyes$ by about 0.17 on the 11B
        (\cref{sec:design-faults}).
  \item Interpretation of low $\capm$: mass has moved to other first tokens,
        typically ``I'' (``I cannot determine\ldots''), measured at 41\% for
        one 11B prompt (diagnostic job 27103565). $\capm$ is therefore both a
        validity statistic and an outcome (\cref{sec:s1-quality}).
  \item State the corollary: $\pyes$ values from conditions with very
        different $\capm$ are conditioned on different fractions of behaviour
        and are not directly comparable.
\end{itemize}
\end{plan}

% ---------------------------------------------------------------------
\section{Item scores and the intrinsic index}\label{sec:framework-asi}
\begin{definition}[Directional bias score and $\ASIint$]\label{def:asi}
For item $i$ with keying $p_i\in\{+1,-1\}$, structure $t$ with sign
$s_t=-1$ for linguistic inversion and $+1$ otherwise, and condition $c$,
\begin{align}
  b_{i,t,c} &= s_t\, p_i\, \pyes^{(i,t,c)}, \label{eq:bias}\\
  \ASIint &= \frac{1}{|I|\,|T|\,|C|}\sum_{i=1}^{22}\sum_{t=1}^{5}\sum_{c=1}^{3} b_{i,t,c}.
  \label{eq:asi}
\end{align}
\end{definition}

\begin{plan}
Explain $s_t$ (``would it be incorrect to say \ldots?'' reverses the meaning of
\emph{yes}); mention that the proposal's Eq.~3.2 omitted $s_t$ and that the
omission was a real bug in the first runner (smoke $\ASIint = 0.639$, fixed).
Give the marginal versions used in Chapter~5 (by structure, condition,
subscale) as one line: replace the triple sum by a sum over the subset.
\end{plan}

% ---------------------------------------------------------------------
\section{Response style, keying imbalance and polarity separation}
\label{sec:framework-responsebias}

\begin{proposition}[Response-style decomposition of $\ASIint$]\label{prop:shift}
Let $\bar p$ be the mean of $\pyes$ over all $22\times5\times3$ cells and
$e_{i,t,c}=\pyes^{(i,t,c)}-\bar p$. Then
\begin{equation}
  \ASIint = \underbrace{\Big(\tfrac{1}{|T|}\textstyle\sum_t s_t\Big)
            \Big(\tfrac{1}{|I|}\textstyle\sum_i p_i\Big)}_{=\;\frac{3}{5}\cdot\frac{10}{22}\;=\;\frac{3}{11}}\,\bar p
            \;+\; \underbrace{\frac{1}{|I||T||C|}\sum_{i,t,c} s_t\,p_i\,e_{i,t,c}}_{D}.
  \label{eq:decomp}
\end{equation}
In particular a model that answers \emph{yes} with the same probability
$\bar p$ to every prompt scores $\ASIint=\tfrac{3}{11}\bar p$, and a uniform
shift $\pyes\mapsto\pyes+a$ changes $\ASIint$ by $\tfrac{3}{11}a$.
\end{proposition}
\begin{proof}
Substitute $\pyes=\bar p+e$ into \cref{eq:asi}. The design is fully crossed, so
the mean of $s_t p_i$ factorises into the mean of $s_t$ over structures
($(4-1)/5$) times the mean of $p_i$ over items ($(16-6)/22$).
\end{proof}

\begin{plan}[Why this matters --- and numbers to quote]
\begin{itemize}
  \item The zero point of $\ASIint$ is not ``no attitude'': an indifferent
        model ($\bar p=0.5$) scores $0.136$.
  \item \newanalysis{(A9) apply \cref{eq:decomp} to every model. Preliminary
        (from the ledger means): Llama-90B $\tfrac{3}{11}\cdot0.5851=0.160$
        vs observed $0.153$, so $D\approx-0.007$; Qwen-72B
        $0.068$ vs $0.008$, $D\approx-0.061$; Llama-11B $0.144$ vs $0.084$,
        $D\approx-0.060$. Llama-90B's headline score is almost entirely its
        yes-rate. Why run it: it turns the thesis-plan's verbal argument into
        an exact identity.}
  \item Link to the psychometric literature on acquiescence with unbalanced
        keying \cite{baumgartner2001response,weijters2013reversed}.
\end{itemize}
\end{plan}

\begin{definition}[Polarity separation]\label{def:ps}
On the non-inversion rows ($t\neq$ inversion), where \emph{yes} always means
agreement with the item as written,
\begin{equation}
  \PS = \overline{\pyes}\big|_{p_i=+1} \;-\; \overline{\pyes}\big|_{p_i=-1},
  \qquad
  \PS_{\logit} = \overline{\logit\pyes}\big|_{p_i=+1} - \overline{\logit\pyes}\big|_{p_i=-1}.
  \label{eq:ps}
\end{equation}
$\PS>0$ means more agreement with sexist than with anti-sexist statements.
\end{definition}

\begin{proposition}[Invariance]\label{prop:invariance}
$\PS$ is unchanged by a uniform shift $\pyes\mapsto\pyes+a$, and $\PS_{\logit}$
is unchanged by a uniform log-odds offset $\logit\pyes\mapsto\logit\pyes+\beta$
(a constant yes-token prior in the sense of \cite{zhao2021calibrate}).
\end{proposition}
\begin{proof}
Both are differences of two means over the same rows; the constant cancels.
\end{proof}

\begin{plan}
\begin{itemize}
  \item \decision{(D3) the August draft used two different definitions
        (all rows in its Table 5.1; non-inversion rows in its \S3.6 and
        Table 5.4). Recommended: \cref{eq:ps} on non-inversion rows as the
        primary metric, $\PS_{\logit}$ and the all-rows version as
        sensitivity checks. With this choice Qwen/Llama is about
        $0.188/0.097\approx1.9\times$, not the draft's $2.3\times$.}
  \item Inversion rows are excluded from $\PS$ because agreement flips there;
        optionally report a ``balanced-keyed'' version that scores inversion
        rows as $1-\pyes$.
\end{itemize}
\end{plan}

% ---------------------------------------------------------------------
\section{Exact permutation inference for polarity}\label{sec:framework-perm}
\begin{plan}
Planned final content: the test statistic and exact $p$-value below. Explain
exchangeability: under $H_0$ the keying labels carry no information about the
model's item-level agreement, so every assignment of the 6 reverse-keyed labels
to the 22 items is equally likely \cite{ernst2004permutation}. The full null
distribution has $\binom{22}{6}=74{,}613$ points, so the test is exact (no
Monte Carlo error). This is the ``permutation test'' promised in proposal
\S3.2.5.
\end{plan}
\begin{equation}
  T(\bm\ell) = \frac{1}{16}\sum_{i:\ell_i=+1}\bar p_i - \frac{1}{6}\sum_{i:\ell_i=-1}\bar p_i,
  \qquad
  p = \frac{1}{\binom{22}{6}}\sum_{\bm\ell\in\Lambda}\ind{|T(\bm\ell)|\ge|T(\bm\ell_{\text{obs}})|},
  \label{eq:perm}
\end{equation}
where $\bar p_i$ is item $i$'s mean $\pyes$ over the 12 non-inversion cells and
$\Lambda$ is the set of all labellings with six reverse-keyed items.

% ---------------------------------------------------------------------
\section{Ratings on a seven-point scale}\label{sec:framework-ratings}
\begin{align}
  \capm_{\Dset} &= \sum_{d=1}^{7} p(d\mid x,v), \qquad
  \rexp = \frac{1}{\capm_{\Dset}}\sum_{d=1}^{7} d\; p(d\mid x,v), \label{eq:rexp}\\
  \rnorm &= \begin{cases} \rexp & \text{original order (7 = extremely objective)}\\
                          8-\rexp & \text{reversed order (7 = extremely subjective)}\end{cases}
  \label{eq:rnorm}
\end{align}
\begin{plan}
Explain why the expectation rather than the argmax (continuous, uses the whole
digit distribution; argmax is stored as a check). The digits are single tokens
only without a leading space, which is why ``Rating:\,'' is prefilled with a
trailing space in the assistant turn (\cref{sec:design-impl}).
\end{plan}

% ---------------------------------------------------------------------
\section{Extrinsic deviation from human baselines}\label{sec:framework-extrinsic}
\begin{equation}
  \dm = \frac{1}{|Q|\,k}\sum_{q\in Q}\sum_{j=1}^{k}\Big(\rnorm^{\,\text{orig}}_{m,q,j}-\rlow\Big),
  \qquad |Q|=10,\; k=9.
  \label{eq:deltam}
\end{equation}

\begin{proposition}[Anchor-insensitivity null]\label{prop:null}
If a model ignores the stated anchors, so that $\rexp^{\,\text{rev}}=\rexp^{\,\text{orig}}$,
then the order-averaged normalised rating is exactly 4 and the order-averaged
deviation for query $q$ is $4-\rlow$. Averaged over the ten SIGIR queries this
null equals $4-5.280=-1.280$.
\end{proposition}
\begin{proof}
$\tfrac12\big(\rexp + (8-\rexp)\big)=4$ for every $\rexp$.
\end{proof}

\begin{equation}
  \lambda_{m,q} = \frac{\bar{\rnorm}_{m,q}-\rlow}{\rhigh-\rlow}
  \qquad (\lambda=0:\ \text{like low-ASI humans};\ \lambda=1:\ \text{like high-ASI humans})
  \label{eq:lambda}
\end{equation}

\begin{plan}
\begin{itemize}
  \item \decision{(D1) which $\dm$. \cref{eq:deltam} uses original-order
        ratings only, as in proposal Eq.~3.4 (one rating per image). The
        pipeline currently averages both CALM orders, which mixes in the
        anchor-insensitivity of \cref{prop:null}: Llama-90B's order-averaged
        $-1.203$ is close to the null $-1.280$. Recommended: original order as
        primary, order-averaged and null-adjusted ($\dm-\dm^{\text{null}}$) as
        sensitivity.}
  \item \decision{(D2) the sign. The proposal reads $\dm<0$ as ``greater
        bias''. But in \cite{otterbacher2018sigir} \emph{more} sexist
        participants were \emph{less} likely to rate biased grids as
        non-objective, i.e.\ they rate them \emph{higher}. A model rating
        below the low-ASI baseline is therefore more ``bias-perceiving'' than
        the least sexist humans, not more sexist. Agree the interpretation
        with your supervisor (a co-author of the SIGIR paper) before writing
        \cref{ch:study2}.}
  \item \cref{eq:lambda} places a model on the low--high ASI axis and is the
        most direct operationalisation of RQ2. \newanalysis{(A10) needs the
        high-ASI baselines $\rhigh$, which require the raw SIGIR CSV (released
        publicly by the SIGIR authors). Guard against small denominators
        $\rhigh-\rlow$.}
\end{itemize}
\end{plan}

% ---------------------------------------------------------------------
\section{Robustness under scale reversal (CALM)}\label{sec:framework-calm}
\begin{align}
  \RR_c(\tau) &= \frac{1}{|D_c|}\sum_{(q,j)\in D_c}
                 \ind{\big|\rnorm^{\,\text{orig}}_{q,j,c}-\rnorm^{\,\text{rev}}_{q,j,c}\big|\le\tau},
  \label{eq:rr}\\
  \dRR(\tau) &= \RR_{\text{vision}}(\tau)-\RR_{\text{text}}(\tau). \label{eq:drr}
\end{align}
\begin{plan}
\begin{itemize}
  \item $\tau=0.5$ (both ratings round to the same scale point). Report
        $\tau\in\{0.25,0.5,1\}$ as a sensitivity table.
  \item \textbf{Consistency rate.} Proposal Eq.~3.6 compares against a
        randomised option order. With only two orders a random order coincides
        with one of them, so $\CR$ is not separately identifiable from $\RR$ in
        this design --- say so explicitly rather than silently dropping it.
  \item \textbf{Effective sample size.} With no image the input is fixed by
        $(q,\text{order})$: the 90 text-only pairs are 10 distinct pairs
        repeated 9 times, so $\RR_{\text{text}}$ is a mean of 10 indicators.
        Inference therefore resamples \emph{queries}: the cluster bootstrap
        \cite{field2007bootstrapping,efron1993bootstrap} draws $B=10{,}000$
        resamples of the 10 queries and reports percentile intervals for
        $\dRR$.
  \item \textbf{Anchor following.} Define the slope $\beta_1$ of
        $\rexp^{\,\text{rev}}$ on $\rexp^{\,\text{orig}}$ across pairs:
        $\beta_1=-1$ means the model follows the reversed anchors,
        $\beta_1=+1$ means it ignores them. This explains \emph{why} $\RR$
        collapses (\cref{sec:s2-calm}).
\end{itemize}
\end{plan}

% ---------------------------------------------------------------------
\section{Validity statistics and degeneracy}\label{sec:framework-validity}
\begin{proposition}[Zero hidden state gives the uniform distribution]\label{prop:uniform}
With $\RMSNorm(\bm h)=\bm g\odot\bm h/\sqrt{\tfrac1d\lVert\bm h\rVert_2^2+\epsilon}$
\cite{zhang2019rmsnorm} and a bias-free output projection, $\bm h_L=\bm 0$
implies $\bm z=\bm 0$ and $p(t\mid x,v)=1/|\vocab|$ for every token. For
Llama-3.2, $1/128{,}256=7.7969\times10^{-6}$ --- exactly the value observed in
the invalid 90B runs.
\end{proposition}
\begin{proof}
$\RMSNorm(\bm 0)=\bm 0$ because $\epsilon>0$; then $\bm z=W_{\text{out}}\bm 0=\bm 0$
and \cref{eq:nexttoken} gives equal probabilities.
\end{proof}

\begin{plan}
Three run-level checks, stated as rules: (i) logit standard deviation at a
fixed self-test prompt must exceed $10^{-4}$; (ii) output diversity
$u=|\{\text{distinct }\pyes\}|/n$ must be near 1 (the invalid 90B run: over
330 prompts, $\pyes$ took 4 distinct values, was exactly 0.5 in 247 rows and
NaN in 59; $\capm$ took 8 distinct non-missing values --- the August draft's
``8 distinct values'' refers to $\capm$, so say which quantity you count);
(iii) report mean and minimum $\capm$ per condition. Point to
\cref{app:fault} for the full fault report.
\end{plan}
